\section{A short criterion for six-port prism blockers}
\label{sec:random-seven-prism-blockers}

This criterion simplifies a necessary condition on critical seven-cut profiles.
It does not establish the unrestricted seven-cut bound.

\begin{lemma}[Prism-blocker criterion]
\label{lem:random-seven-prism-blocker}
Let $D$ be a graph on a fixed set of six vertices.
The following conditions are equivalent.
\begin{enumerate}
\item There is a six-cycle $C$ in $K_6$ such that $D$ meets every perfect matching of $K_6-E(C)$.
\item Every partition of the six vertices into two triples has an edge of $D$ between its parts.
\item The graph $D$ contains a copy of $K_{1,3}$, $P_4$, or three disjoint edges.
\end{enumerate}
\end{lemma}

\begin{proof}
First suppose that $D$ has no edge across a partition into triples $A,B$.
Fix any six-cycle $C$.
The bipartite graph $K_{A,B}-E(C)$ has minimum degree at least one.
If two vertices on one side had only one common available neighbor, the four deleted cross edges would form a four-cycle in $C$.
This is impossible.
Hall's condition therefore holds for subsets of sizes one, two, and three.
There is a perfect matching across $A,B$ that avoids $C$.
It also avoids $D$.
This proves that the first condition implies the second.

If a component of $D$ has at least four vertices, it contains a tree on four vertices.
Such a tree is $K_{1,3}$ or $P_4$.
Otherwise all components have at most three vertices.
A component of order three immediately gives a partition into triples with no crossing edge.
When all components have order at most two, such a partition exists unless their orders are $2,2,2$.
The remaining case contains three disjoint edges.
This proves that the second condition implies the third.

Every triangular prism is the complement of a six-cycle in $K_6$.
Its perfect matchings consist of all three spokes, or one spoke together with one edge from each triangle.
Three disjoint edges meet all four matchings when used as the three spokes.
A three-edge star meets all perfect matchings when used as the edges incident with one prism vertex.

For a path $0\,1\,2\,3$, take prism triangles $\{0,1,5\}$ and $\{2,3,4\}$, with spokes $04,12,35$.
The all-spoke matching and the matching with single spoke $12$ meet the path in edge $12$.
The matching with single spoke $04$ contains edge $23$.
The matching with single spoke $35$ contains edge $01$.
Thus all three listed subgraphs give the first condition.
\end{proof}

\begin{corollary}[Ten inequalities for an active six-sector]
\label{cor:random-seven-balanced-prism}
Consider a critical seven-cut profile with common Boolean sector totals $z_Q$.
Let $S$ be a six-port set and let $S=A\sqcup B$, where $|A|=|B|=3$.
Then
\[
 \sum_{a\in A,\ b\in B}\bigl(1-z_{S\setminus\{a,b\}}\bigr)
 \ge z_S.
\]
Equivalently,
\[
 \sum_{a\in A,\ b\in B}z_{S\setminus\{a,b\}}+z_S\le9.
\]
There are ten such inequalities for each six-sector.
\end{corollary}

\begin{proof}
If $z_S=0$, the inequality follows from the Boolean bounds.
Suppose that $z_S=1$.
Choose positive six-terminal pairing masses on the two shores.
Their matchings $\beta,\gamma$ are disjoint, so their union is a six-cycle.
Each common compatible matching is a perfect matching of $K_6-(\beta\cup\gamma)$.
The matching-augmentation constraints and disjoint four-pairing supports require every such matching to contain a pair whose complementary four-sector is inactive.
Apply Lemma~\ref{lem:random-seven-prism-blocker} to these inactive complementary pairs.
\end{proof}

The criterion also shows that the maximal missing-pair sets failing this necessary condition are exactly the ten unions of two disjoint triangles.
The compact numerical relaxation uses these seventy inequalities together with the local star-cover bounds.
Numerical feasibility or infeasibility in that relaxation remains separate from a graph theorem.
